\documentclass{beamer}
\usetheme{Madrid}
\usepackage[utf8]{inputenc}
\usepackage{amsmath}
\usepackage{amsfonts}
\usepackage{amssymb}
\usepackage{graphicx}
\usepackage{xcolor}
\usepackage{tikz}
\usepackage{booktabs}
\usepackage{array}
\usepackage{colortbl}
\usepackage{multicol}

% ============================================================
% TINY FONT FOR MAXIMUM CONTENT
% ============================================================
\usefonttheme{professionalfonts}
\setbeamerfont{normal text}{size=\scriptsize}
\setbeamerfont{itemize/enumerate body}{size=\scriptsize}
\setbeamerfont{itemize/enumerate subbody}{size=\tiny}
\setbeamerfont{block body}{size=\scriptsize}
\setbeamerfont{block title}{size=\scriptsize}
\setbeamerfont{frametitle}{size=\small}
\setbeamerfont{title}{size=\normalsize}
\setbeamerfont{subtitle}{size=\small}
\setbeamerfont{author}{size=\tiny}
\setbeamerfont{date}{size=\tiny}

\renewcommand{\scriptsize}{\fontsize{6pt}{7pt}\selectfont}
\renewcommand{\footnotesize}{\fontsize{6.5pt}{8pt}\selectfont}
\renewcommand{\small}{\fontsize{7.5pt}{9pt}\selectfont}
\renewcommand{\normalsize}{\fontsize{8.5pt}{10.5pt}\selectfont}

\newcommand{\redflag}[1]{\textcolor{red}{\textbf{[!] #1}}}
\newcommand{\bluenote}[1]{\textcolor{blue}{\textbf{[i] #1}}}
\newcommand{\greennote}[1]{\textcolor{green!50!black}{\textbf{[Right] #1}}}
\newcommand{\examfocus}[1]{\textcolor{orange!80!black}{\textbf{[FOCUS] #1}}}
\newcommand{\trap}[1]{\textcolor{purple}{\textbf{[TRAP] #1}}}
\newcommand{\correct}[1]{\textcolor{green!50!black}{\textbf{[CORRECT] #1}}}
\newcommand{\scenario}[1]{\textcolor{blue!70!black}{\textbf{[SCENARIO] #1}}}
\newcommand{\keyterm}[1]{\textcolor{red!80!black}{\textbf{[KEY] #1}}}
\newcommand{\lawcite}[1]{\textcolor{brown!70!black}{\textbf{[LAW] #1}}}

\title[CAMS Domain A]{CAMS Exam Preparation\\Domain A: Understanding the Risks and Methods of Financial Crime (30\%)}
\author{Advanced AML Training Framework}
\date{}

\setbeamercovered{transparent}
\setbeamertemplate{navigation symbols}{}
\setbeamertemplate{frametitle continuation}{}
\setlength{\parskip}{2pt}
\setlength{\itemsep}{1pt}

\begin{document}
	
	% ============================================================
	% TITLE SLIDE
	% ============================================================
	\begin{frame}
		\titlepage
		\centering
		\tiny Domain A: 30\% of CAMS Exam | 100 Slides | ML Stages, TBML, PEPs, VASPs, Gatekeepers, Real Estate, Insurance
	\end{frame}
	
	% ============================================================
	% SECTION 1: DEFINITIONS OF AML, CFT, SANCTIONS, FRAUD, ABC, TAX EVASION
	% ============================================================
	
	\begin{frame}{1.1: Money Laundering vs. Terrorist Financing}
		\begin{block}{Key Distinction (Heavily Tested)}
			\begin{itemize}
				\item \textbf{Money Laundering}: Conceals \textbf{illicit origin} of proceeds; always involves dirty money from predicate crimes.
				\item \textbf{Terrorist Financing}: Funds may be \textbf{legitimate} (donations, salary); concealment objective is \textbf{intended use} for violence, not origin.
			\end{itemize}
		\end{block}
		\textcolor{red}{\textbf{CRITICAL:}} Standard AML controls focused on suspicious \textbf{origins} may fail to detect TF involving \textbf{clean-source funds}.
	\end{frame}
	
	\begin{frame}{1.2: Money Laundering Definition}
		\begin{block}{Core Definition}
			The process of making criminal proceeds appear legitimate through three distinct stages: Placement, Layering, Integration.
		\end{block}
		\textbf{Key Elements:}
		\begin{itemize}
			\item Criminal proceeds from predicate offenses
			\item Intent to conceal illicit origin
			\item Three-stage process
			\item Goal: Make funds appear legitimate
		\end{itemize}
	\end{frame}
	
	\begin{frame}{1.3: Terrorist Financing Definition}
		\begin{block}{FATF Recommendation 5 - TF Offense}
			Criminalize TF regardless of whether funds actually used for terrorist act.
		\end{block}
		\textbf{Requirements:}
		\begin{itemize}
			\item No minimum monetary threshold (even \$1 is enough)
			\item Collection of funds for terrorist purposes is sufficient
			\item Standalone offense (not just covered under ML)
			\item Intent is the critical element
		\end{itemize}
		\textcolor{red}{\textbf{Deficiency:}} \$10,000 threshold for TF violates Recommendation 5.
	\end{frame}
	
	\begin{frame}{1.4: Sanctions Definition}
		\begin{block}{What Are Sanctions?}
			Legal restrictions imposed by governments or international bodies against designated countries, individuals, or entities to achieve foreign policy or national security objectives.
		\end{block}
		\textbf{Types of Sanctions:}
		\begin{itemize}
			\item Country/Regime Sanctions: Comprehensive restrictions
			\item Targeted Sanctions: Against specific individuals or entities
			\item Sectoral Sanctions: Restrictions on specific industries
			\item Asset Freezes: Freezing assets of designated persons
		\end{itemize}
		\textcolor{red}{Sanctions are \textbf{strict liability} — no intent is required for a violation.}
	\end{frame}
	
	\begin{frame}{1.5: Fraud Definition}
		\begin{block}{Financial Fraud Definition}
			Intentional deception or misrepresentation made for personal gain or to damage another individual, involving financial transactions or assets.
		\end{block}
		\textbf{Common Fraud Types:}
		\begin{itemize}
			\item Identity Theft: Using stolen personal information
			\item Synthetic Identity Fraud: Creating fake identities
			\item Investment Fraud: False investment opportunities
			\item Payment Fraud: Unauthorized transactions
			\item Insurance Fraud: False insurance claims
		\end{itemize}
		\textcolor{blue}{Fraud and money laundering often overlap — fraud generates proceeds that need laundering.}
	\end{frame}
	
	\begin{frame}{1.6: Anti-Bribery and Corruption Definition}
		\begin{block}{Bribery Definition}
			The offering, giving, receiving, or soliciting of anything of value to influence an action or decision.
		\end{block}
		\textbf{Key Distinctions:}
		\begin{itemize}
			\item \textbf{Bribery}: Paying for influence (illegal)
			\item \textbf{Facilitation Payment}: Small payment to speed routine action (allowed under FCPA, prohibited under UK Bribery Act)
			\item \textbf{Hospitality}: Legitimate business entertainment (must be proportionate)
			\item \textbf{Corruption}: Abuse of entrusted power for private gain
		\end{itemize}
		\textcolor{red}{The \textbf{UK Bribery Act} is stricter than the FCPA — it prohibits facilitation payments.}
	\end{frame}
	
	\begin{frame}{1.7: Tax Evasion Definition}
		\begin{block}{Tax Evasion Definition}
			The illegal act of deliberately not paying taxes that are legally owed, typically through concealment or misrepresentation.
		\end{block}
		\textbf{Distinction:}
		\begin{itemize}
			\item \textbf{Tax Evasion}: Illegal — misrepresentation, concealment
			\item \textbf{Tax Avoidance}: Legal — using legitimate methods to reduce tax liability
		\end{itemize}
		\textcolor{red}{Tax evasion is a predicate offense for money laundering in many jurisdictions.}
	\end{frame}
	
	\begin{frame}{1.8: FATF Recommendation 3 - Predicate Offenses}
		\begin{block}{Key Requirements}
			\begin{itemize}
				\item FATF recommends the \textbf{"all-crimes approach"} - applies to proceeds from \textbf{all serious criminal acts}.
				\item \textbf{Self-laundering IS criminalized} - predicate offender can also be charged with money laundering.
				\item 6AMLD harmonizes \textbf{22 predicate offenses} across EU.
			\end{itemize}
		\end{block}
		\textcolor{blue}{Example:} Drug trafficker who processes own proceeds through a car wash business $\rightarrow$ can be charged with BOTH drug trafficking AND money laundering.
	\end{frame}
	
	% ============================================================
	% SECTION 2: FINANCIAL SYSTEMS AND CONSEQUENCES
	% ============================================================
	
	\begin{frame}{2.1: Economic Consequences of ML}
		\begin{block}{Economic Impacts}
			\begin{itemize}
				\item Distortion of economic statistics and misallocation of resources
				\item Undermines legitimate economic activity
				\item Loss of tax revenue for governments
				\item Increased volatility from sudden capital movements
				\item Reduced foreign investment in high-risk countries
				\item Reputational damage to financial centers
			\end{itemize}
		\end{block}
		\textcolor{red}{Money laundering \textbf{undermines} legitimate economic activity by distorting markets.}
	\end{frame}
	
	\begin{frame}{2.2: Reputational Consequences of ML}
		\begin{block}{Reputational Impacts}
			\begin{itemize}
				\item Loss of customer trust and confidence
				\item Reduced business opportunities and partnerships
				\item Damage to brand value and market position
				\item Negative media coverage and public scrutiny
				\item Increased regulatory supervision and oversight
				\item Difficulty attracting and retaining talent
			\end{itemize}
		\end{block}
		\textcolor{red}{Reputational damage can be \textbf{more severe} than financial penalties.}
	\end{frame}
	
	\begin{frame}{2.3: Social Consequences of ML}
		\begin{block}{Social Impacts}
			\begin{itemize}
				\item Increased crime as laundered funds fuel criminal enterprises
				\item Corruption of public officials and institutions
				\item Social inequality from illicit wealth concentration
				\item Weakens rule of law and legal institutions
				\item Fuels violence, terrorism, and organized crime
				\item Increases poverty and undermines development
			\end{itemize}
		\end{block}
		\textcolor{red}{Money laundering is not a victimless crime — it fuels \textbf{organized crime} and \textbf{corruption}.}
	\end{frame}
	
	\begin{frame}{2.4: Regulatory Penalties for AML Violations}
		\begin{block}{Types of Penalties}
			\begin{itemize}
				\item \textbf{Civil Penalties}: Financial fines (e.g., US BSA, UK FCA)
				\item \textbf{Criminal Penalties}: Imprisonment for individuals
				\item \textbf{Administrative Penalties}: License suspension or revocation
				\item \textbf{Personal Liability}: Compliance officers and senior management
			\end{itemize}
		\end{block}
		\textcolor{red}{Under \textbf{FATF Recommendation 35}, sanctions must be effective, proportionate, and dissuasive.}
	\end{frame}
	
	\begin{frame}{2.5: Consequences of Indiscriminate De-Risking}
		\begin{block}{De-Risking Definition}
			Termination of entire categories of customer relationships without individual risk assessment.
		\end{block}
		\textbf{Negative Consequences:}
		\begin{itemize}
			\item Legitimate businesses lose access to global financial system
			\item Drives transactions into \textbf{unregulated, informal channels}
			\item Harms financial inclusion and developing economies
			\item Contradicts the mandatory \textbf{risk-based approach}
		\end{itemize}
	\end{frame}
	
	% ============================================================
	% SECTION 3: PREDICATE CRIMES
	% ============================================================
	
	\begin{frame}{3.1: Predicate Offenses Overview}
		\begin{block}{FATF Predicate Offense Categories}
			\begin{itemize}
				\item UN Convention Categories: Drug trafficking, corruption, fraud, human trafficking
				\item FATF Expansion: Tax crimes, environmental crimes, counterfeiting
				\item EU 6AMLD: Expanded to include 22 predicate offenses
			\end{itemize}
		\end{block}
		\textcolor{red}{Under FATF standards, \textbf{all serious offenses} should be predicate offenses for ML.}
	\end{frame}
	
	\begin{frame}{3.2: Drug Trafficking and ML}
		\begin{block}{Drug Trafficking as ML Source}
			Drug trafficking generates enormous cash proceeds that must be laundered to be used in the legitimate economy.
		\end{block}
		\textbf{Key Points:}
		\begin{itemize}
			\item Cash-intensive nature makes it a major ML source
			\item Proceeds often laundered through trade-based ML
			\item Cartels use BMPE and trade to move value
			\item Drug trafficking organizations use complex structures
		\end{itemize}
	\end{frame}
	
	\begin{frame}{3.3: Corruption and ML}
		\begin{block}{Corruption-ML Connection}
			Corruption proceeds (bribes, embezzlement) must be laundered through the financial system to appear legitimate.
		\end{block}
		\textbf{Key Points:}
		\begin{itemize}
			\item PEPs are at higher risk of corruption
			\item Corruption enables other financial crimes
			\item Bribery undermines internal controls
			\item Corruption and money laundering are \textbf{interconnected}
		\end{itemize}
	\end{frame}
	
	\begin{frame}{3.4: Human Trafficking and ML}
		\begin{block}{Human Trafficking-ML Connection}
			Human trafficking generates significant proceeds that must be laundered through the financial system.
		\end{block}
		\textbf{Key Points:}
		\begin{itemize}
			\item High-profit, low-risk crime for traffickers
			\item Multiple small payments to controller accounts
			\item Funds used to purchase real estate and businesses
			\item Victims' wages deposited to controller accounts
		\end{itemize}
	\end{frame}
	
	\begin{frame}{3.5: Self-Laundering and Double Jeopardy}
		\begin{block}{Self-Laundering (18 U.S.C. \$1956)}
			U.S. law \textbf{DOES} criminalize self-laundering.
		\end{block}
		\textbf{Key Points:}
		\begin{itemize}
			\item Predicate offender CAN be charged with money laundering for processing own proceeds
			\item \textbf{Double jeopardy argument fails} - each offense requires proof of different elements
			\item Example: Wire fraud + money laundering = TWO separate crimes
		\end{itemize}
	\end{frame}
	
	% ============================================================
	% SECTION 4: HOW FINANCIAL CRIME MANIFESTS
	% ============================================================
	
	\begin{frame}{4.1: The Three Stages of Money Laundering}
		\begin{block}{Three Stages}
			\begin{enumerate}
				\item \textbf{Placement}: Introducing illicit funds into the financial system
				\item \textbf{Layering}: Conducting complex transactions to obscure the audit trail
				\item \textbf{Integration}: Making laundered funds appear legitimate through investments
			\end{enumerate}
		\end{block}
		\textcolor{blue}{Example:} Cash deposits $\rightarrow$ wire transfers $\rightarrow$ real estate purchase.
	\end{frame}
	
	\begin{frame}{4.2: Placement Techniques}
		\begin{block}{Common Placement Methods}
			\begin{itemize}
				\item Structured cash deposits (smurfing) below CTR thresholds
				\item Currency smuggling across borders
				\item Purchasing assets with cash
				\item Using currency exchanges or MSBs
				\item Gambling at casinos
				\item Commingling with legitimate business revenues
			\end{itemize}
		\end{block}
		\textcolor{red}{Placement is the riskiest stage for criminals.}
	\end{frame}
	
	\begin{frame}{4.3: Structuring / Smurfing}
		\begin{block}{Structuring Definition}
			Deliberately breaking cash deposits into amounts below reporting thresholds to evade CTR requirements.
		\end{block}
		\textbf{Red flags:}
		\begin{itemize}
			\item Multiple deposits just below \$10,000 (\$9,500-\$9,900)
			\item Deposits from multiple individuals to same account
			\item Geographic dispersion of deposits
		\end{itemize}
		\textcolor{red}{Structuring is a federal crime independent of the underlying source of funds.}
	\end{frame}
	
	\begin{frame}{4.4: Layering Techniques}
		\begin{block}{Layering Methods}
			\begin{itemize}
				\item Wire transfers through multiple accounts and jurisdictions
				\item Use of shell companies and trusts
				\item Trade-based money laundering
				\item Cryptocurrency mixing and privacy coin conversion
				\item Use of professional intermediaries
			\end{itemize}
		\end{block}
		\textbf{Purpose:} Create distance between funds and criminal source.
	\end{frame}
	
	\begin{frame}{4.5: Integration Techniques}
		\begin{block}{Integration Methods}
			\begin{itemize}
				\item Purchase of real estate or luxury assets
				\item Investment in legitimate businesses
				\item False loan repayments
				\item Commingling with legitimate business revenues
				\item Sale of assets purchased with laundered funds
			\end{itemize}
		\end{block}
		\textbf{Purpose:} Funds re-enter economy as legitimate-appearing wealth.
	\end{frame}
	
	\begin{frame}{4.6: Over-Invoicing (TBML)}
		\begin{block}{Over-Invoicing Mechanism}
			\begin{itemize}
				\item Exporter inflates invoice price above market value
				\item Importer pays inflated amount, transferring excess value to exporter
			\end{itemize}
		\end{block}
		\textbf{Red flag:} Price significantly above market (e.g., \$1,200 vs \$180).
		\textcolor{blue}{Criminal purpose:} Move funds from importer to exporter disguised as trade.
	\end{frame}
	
	\begin{frame}{4.7: Under-Invoicing (TBML)}
		\begin{block}{Under-Invoicing Mechanism}
			\begin{itemize}
				\item Exporter deflates invoice price below market value
				\item Importer pays lower amount but sells goods at market price
				\item Importer retains surplus as "clean" funds
			\end{itemize}
		\end{block}
		\textbf{Red flag:} Price significantly below market.
		\textcolor{blue}{Criminal purpose:} Move value from exporter to importer.
	\end{frame}
	
	\begin{frame}{4.8: Phantom Shipping}
		\begin{block}{Phantom Shipping Mechanism}
			\begin{itemize}
				\item No actual goods are ever shipped
				\item Complete fabricated documentation
				\item Funds transferred based on false paperwork
			\end{itemize}
		\end{block}
		\textbf{Detection indicators:}
		\begin{itemize}
			\item Empty containers discovered upon inspection
			\item Shipper has no manufacturing capability
			\item Altered bills of lading
		\end{itemize}
	\end{frame}
	
	\begin{frame}{4.9: Free Trade Zone (FTZ) Vulnerabilities}
		\begin{block}{Why FTZs are High Risk}
			\begin{itemize}
				\item Relaxed regulatory oversight
				\item Minimal customs documentation inspections
				\item High volume of re-exported goods
				\item Opportunity for mislabeling cargo
				\item Phantom shipments
			\end{itemize}
		\end{block}
		\textbf{Red flags:} Uniform invoice values, empty containers, altered bills of lading.
	\end{frame}
	
	\begin{frame}{4.10: Black Market Peso Exchange (BMPE)}
		\begin{block}{BMPE Mechanism}
			\begin{enumerate}
				\item Narco-dollars deposited in U.S. accounts (structured below \$10k)
				\item Colombian peso broker buys dollars at discount (20-30\% off)
				\item Broker uses dollars to pay U.S. invoices for Colombian merchants
				\item Colombian merchant repays broker in pesos at favorable rate
				\item \textbf{No cross-border currency movement} after initial deposit
			\end{enumerate}
		\end{block}
		\textbf{Key insight:} Uses drug dollars to pay legitimate commercial invoices.
	\end{frame}
	
	% ============================================================
	% SECTION 5: BANKING PRODUCT RISKS
	% ============================================================
	
	\begin{frame}{5.1: Trade Finance Risks}
		\begin{block}{Trade Finance ML Risks}
			Over/under-invoicing to move value across borders while disguising the true purpose of the transaction.
		\end{block}
		\textbf{Red Flags:}
		\begin{itemize}
			\item Invoice price significantly above or below market value
			\item Multiple amendments to letters of credit
			\item Round-dollar invoices (not typical for genuine trade)
			\item Shipment routing through free trade zones
		\end{itemize}
		\textcolor{red}{Trade-based ML is \textbf{difficult to detect} because it appears as legitimate trade.}
	\end{frame}
	
	\begin{frame}{5.2: Wealth Management Risks}
		\begin{block}{Wealth Management ML Risks}
			Use of offshore structures, trusts, and complex investment vehicles to conceal ownership and source of funds.
		\end{block}
		\textbf{Key Vulnerabilities:}
		\begin{itemize}
			\item Complexity and confidentiality attract launderers
			\item Numbered accounts and relationship discretion
			\item High asset requirements
			\item Use of multiple intermediaries
		\end{itemize}
	\end{frame}
	
	\begin{frame}{5.3: Correspondent Banking Risks}
		\begin{block}{Correspondent Banking Risks}
			The respondent bank may have weaker AML controls, creating a channel for illicit funds to enter the financial system.
		\end{block}
		\textbf{Key Vulnerabilities:}
		\begin{itemize}
			\item Loss of visibility into underlying transactions
			\item Nested correspondent accounts
			\item Respondent bank services unlicensed MSBs
			\item Payable-through accounts with no CDD
		\end{itemize}
		\textcolor{red}{Correspondent banking is a \textbf{gateway} for cross-border money laundering.}
	\end{frame}
	
	\begin{frame}{5.4: Private Banking Risks}
		\begin{block}{Private Banking Risks}
			Discretion and personalized service can be exploited to conceal illicit activity and bypass standard controls.
		\end{block}
		\textbf{Key Vulnerabilities:}
		\begin{itemize}
			\item Numbered accounts
			\item Relationship discretion
			\item High asset requirements
			\item Personalized service can circumvent controls
		\end{itemize}
	\end{frame}
	
	\begin{frame}{5.5: Nested Correspondent Accounts}
		\begin{block}{Nested Correspondent Accounts - Critical Risk}
			When a respondent bank allows third parties to route transactions through the correspondent account \textbf{without disclosure}:
			\begin{itemize}
				\item Correspondent bank loses visibility into true originators/beneficiaries
				\item Increased risk of processing for sanctioned entities
				\item Violates core correspondent banking due diligence
			\end{itemize}
		\end{block}
		\textbf{Required actions:} Immediate RFI, SAR filing, reassess relationship.
	\end{frame}
	
	\begin{frame}{5.6: Payable-Through Accounts (PTAs)}
		\begin{block}{PTA Obligations (U.S. Regulations)}
			\begin{itemize}
				\item Correspondent bank MUST apply CDD directly to underlying PTA accountholders
				\item Must maintain records identifying each PTA accountholder
				\item \textbf{Cannot rely solely on respondent bank's CDD}
			\end{itemize}
		\end{block}
		\textbf{Red flag:} Multiple sub-\$10,000 wires to same shell company.
	\end{frame}
	
	\begin{frame}{5.7: Shell Bank Prohibition (Section 313)}
		\begin{block}{Section 313 - Shell Bank Definition}
			A "shell bank" has:
			\begin{itemize}
				\item NO physical presence in any country
				\item NOT affiliated with a regulated financial institution
			\end{itemize}
		\end{block}
		\textcolor{red}{\textbf{U.S. banks absolutely prohibited}} from maintaining correspondent accounts for foreign shell banks.
	\end{frame}
	
	% ============================================================
	% SECTION 6: BANKING HIGH-RISK SECTORS
	% ============================================================
	
	\begin{frame}{6.1: Banking DNFBPs (Lawyers, Accountants)}
		\begin{block}{DNFBP Risks}
			DNFBPs handle client money and can facilitate transactions that conceal the source of funds.
		\end{block}
		\textbf{Key Vulnerabilities:}
		\begin{itemize}
			\item Gatekeepers who can enable money laundering
			\item Professional status may bypass scrutiny
			\item IOLTA accounts are blind pass-through vehicles
		\end{itemize}
		\textcolor{red}{DNFBPs are \textbf{gatekeepers} who can enable money laundering without knowing it.}
	\end{frame}
	
	\begin{frame}{6.2: Banking Casinos}
		\begin{block}{Casino Risks}
			Casinos handle large amounts of cash and can be used for placement and layering of illicit funds.
		\end{block}
		\textbf{Key Vulnerabilities:}
		\begin{itemize}
			\item Cash-intensive businesses
			\item Can commingle illicit and legitimate funds
			\item Chip washing combines placement and integration
		\end{itemize}
	\end{frame}
	
	\begin{frame}{6.3: Banking Real Estate}
		\begin{block}{Real Estate Risks}
			Real estate purchases can be used to launder large sums of money through asset purchases.
		\end{block}
		\textbf{Key Vulnerabilities:}
		\begin{itemize}
			\item Preferred integration vehicle
			\item All-cash purchases through anonymous LLCs
			\item Over/under-valuation techniques
		\end{itemize}
	\end{frame}
	
	\begin{frame}{6.4: Banking Cryptocurrency}
		\begin{block}{Cryptocurrency Risks}
			Pseudonymous transactions and rapid value transfer make crypto exchanges attractive for money laundering.
		\end{block}
		\textbf{Key Vulnerabilities:}
		\begin{itemize}
			\item Pseudonymity and anonymity
			\item Rapid cross-border transfer
			\item Privacy coins and mixers
		\end{itemize}
		\textcolor{red}{Cryptocurrency exchanges are \textbf{VASPs} and are subject to FATF Recommendation 16 (Travel Rule).}
	\end{frame}
	
	% ============================================================
	% SECTION 7: PEP AND HIGH-RISK CUSTOMER RISKS
	% ============================================================
	
	\begin{frame}{7.1: PEP Definition}
		\begin{block}{Who Qualifies as PEP (FATF Definition)}
			\textbf{Individuals with prominent public functions:}
			\begin{itemize}
				\item Heads of state or government
				\item Senior politicians and cabinet ministers
				\item Senior judicial officials
				\item Senior military officials
				\item Senior executives of state-owned enterprises
				\item Leaders of political parties
			\end{itemize}
			\textbf{Family members and close associates} are also treated as PEPs.
		\end{block}
	\end{frame}
	
	\begin{frame}{7.2: PEP Family Members and Close Associates}
		\begin{block}{Family Member Definition}
			Treated as PEPs even without government positions:
			\begin{itemize}
				\item Spouse or civil partner
				\item Children (including adult children)
				\item Parents
				\item Siblings
			\end{itemize}
			\textbf{Close associates:}
			\begin{itemize}
				\item Close business or personal relationship to PEP
				\item Beneficial owners of entities for PEP's benefit
			\end{itemize}
		\end{block}
		\textcolor{red}{Example:} Adult daughter of foreign minister IS a PEP requiring EDD.
	\end{frame}
	
	\begin{frame}{7.3: Foreign vs. Domestic PEPs}
		\begin{block}{Classification Differences}
			\textbf{Foreign PEP:}
			\begin{itemize}
				\item Always requires EDD and senior approval
				\item Higher inherent risk
				\item Must establish SoW and SoF
			\end{itemize}
			\textbf{Domestic PEP:}
			\begin{itemize}
				\item May require EDD depending on jurisdiction
				\item Potentially less intensive
				\item Still requires enhanced scrutiny
			\end{itemize}
		\end{block}
	\end{frame}
	
	\begin{frame}{7.4: Source of Wealth vs. Source of Funds}
		\begin{block}{Critical Distinction (Heavily Tested)}
			\begin{itemize}
				\item \textbf{Source of Funds (SoF)}: Specific transaction being deposited (sale contract, inheritance)
				\item \textbf{Source of Wealth (SoW)}: How PEP accumulated overall net worth (career, business, inheritance)
			\end{itemize}
		\end{block}
		\textcolor{red}{\textbf{EDD for PEPs requires BOTH.}} Civil servant with \$4.5M apartment: SoF insufficient without SoW.
	\end{frame}
	
	\begin{frame}{7.5: PEP Lifecycle - Mid-Relationship Designation Change}
		\begin{block}{When a Customer Becomes a PEP}
			\begin{itemize}
				\item \textbf{Immediate reclassification to HIGH risk}
				\item Mandatory EDD applied within days
				\item Re-establish source of wealth AND source of funds
				\item Senior management approval to continue
				\item Back-dated gap period is control failure
			\end{itemize}
		\end{block}
		\textcolor{blue}{Example:} Customer elected to foreign cabinet $\rightarrow$ re-rate immediately.
	\end{frame}
	
	\begin{frame}{7.6: Former PEP - Grace Period}
		\begin{block}{PEP Status After Retirement}
			\begin{itemize}
				\item Maintain for \textbf{"reasonable period"} (minimum 12 months)
				\item \textbf{Risk-based decision} considering:
				\begin{itemize}
					\item Country corruption risk
					\item Continued access to influence
					\item Time elapsed since leaving office
					\item Account activity pattern
				\end{itemize}
				\item 5 years clean activity MAY justify lowering risk
				\item \textbf{Decision must be documented}
			\end{itemize}
		\end{block}
	\end{frame}
	
	\begin{frame}{7.7: EDD - Required Actions for PEPs}
		\begin{block}{Enhanced Due Diligence Requirements}
			\begin{itemize}
				\item Establish SoW and SoF through verifiable documentation
				\item Obtain senior management approval
				\item Increased transaction monitoring frequency
				\item Review all significant transactions
				\item Identify UBO of any legal entity or trust
				\item Document all findings
				\item Reassess risk annually or upon triggers
			\end{itemize}
		\end{block}
	\end{frame}
	
	% ============================================================
	% SECTION 8: OTHER FINANCIAL SECTOR PRODUCT RISKS
	% ============================================================
	
	\begin{frame}{8.1: Securities Sector Risks}
		\begin{block}{Securities Sector ML Risks}
			Market manipulation and insider trading generate illicit profits that must be laundered.
		\end{block}
		\textbf{Key Vulnerabilities:}
		\begin{itemize}
			\item Used for both generating and laundering illicit funds
			\item Complex trading patterns obscure source
			\item Multiple accounts and jurisdictions
		\end{itemize}
	\end{frame}
	
	\begin{frame}{8.2: Investment Fund Risks}
		\begin{block}{Investment Fund ML Risks}
			Use of funds to commingle illicit and legitimate investments, obscuring the source of funds.
		\end{block}
		\textbf{Key Vulnerabilities:}
		\begin{itemize}
			\item Integration vehicle for laundered funds
			\item Complex structures hide ownership
			\item Multiple layers of funds
		\end{itemize}
	\end{frame}
	
	\begin{frame}{8.3: Leasing and Financing Risks}
		\begin{block}{Asset Leasing and Financing Risks}
			Over-valuation of assets to move value, or using lease payments to layer funds.
		\end{block}
		\textbf{Key Vulnerabilities:}
		\begin{itemize}
			\item Over-invoicing and value transfer
			\item Complex financing structures
			\item Cross-border leasing arrangements
		\end{itemize}
	\end{frame}
	
	\begin{frame}{8.4: Pension Fund Risks}
		\begin{block}{Pension Fund ML Risks}
			Illicit funds can be invested in pension funds, later withdrawn as "clean" retirement benefits.
		\end{block}
		\textbf{Key Vulnerabilities:}
		\begin{itemize}
			\item Long-term integration vehicle
			\item Legitimate-appearing retirement income
			\item Multiple layers of investments
		\end{itemize}
	\end{frame}
	
	% ============================================================
	% SECTION 9: MSB, PSP, E-COMMERCE RISKS
	% ============================================================
	
	\begin{frame}{9.1: MSB Registration Triggers (U.S. FinCEN)}
		\begin{block}{Who Must Register as an MSB}
			\begin{itemize}
				\item Check casher: cashing > \$1,000/person/day
				\item Money transmitter: any volume
				\item Currency dealer or exchanger
				\item Issuer or seller of prepaid access
				\item Traveler's check issuer
			\end{itemize}
		\end{block}
		\textcolor{red}{Grocery store cashing checks > \$1,000/day without registration = BSA violation.}
	\end{frame}
	
	\begin{frame}{9.2: MSB Principal Responsibility for Agents}
		\begin{block}{Agent Network Compliance}
			\begin{itemize}
				\item MSB principal is \textbf{fully responsible} for agent compliance
				\item Must monitor agent transaction activity
				\item Unexplained volume spikes require investigation
				\item May need to file SAR on agent location
				\item Terminate agent if compliance cannot be ensured
			\end{itemize}
		\end{block}
		\textbf{Red flag:} Agent processes 5,000\% above historical average.
	\end{frame}
	
	\begin{frame}{9.3: PSP Risks}
		\begin{block}{Payment Service Provider Risks}
			Facilitating fraudulent transactions and operating in jurisdictions with limited regulatory oversight.
		\end{block}
		\textbf{Key Vulnerabilities:}
		\begin{itemize}
			\item Gateways for e-commerce fraud
			\item Cross-border payment complexity
			\item Multiple payment methods obscure source
		\end{itemize}
	\end{frame}
	
	\begin{frame}{9.4: E-Commerce Platform Risks}
		\begin{block}{E-Commerce ML Risks}
			Using fake transactions or over/under-invoicing to move value across borders.
		\end{block}
		\textbf{Key Vulnerabilities:}
		\begin{itemize}
			\item Trade-based money laundering
			\item Fake transactions and refunds
			\item Multiple seller accounts
		\end{itemize}
	\end{frame}
	
	\begin{frame}{9.5: Prepaid Access Programs}
		\begin{block}{Prepaid Access Requirements}
			\begin{itemize}
				\item Issuer is considered an MSB
				\item Must register with FinCEN
				\item Must implement risk-based AML program
				\item CDD on cardholders exceeding thresholds
				\item Monitor reload activity for structuring
			\end{itemize}
		\end{block}
		\textbf{Not exempt:} Prepaid cards are monetary instruments under BSA.
	\end{frame}
	
	\begin{frame}{9.6: Prepaid Cards - Structuring and Border Evasion}
		\begin{block}{Prepaid Card Abuse Example}
			Traveler with 50 prepaid cards each loaded with \$9,500 (total \$475,000):
			\begin{itemize}
				\item \textbf{Structuring violation}: Each card just below \$10,000
				\item \textbf{Currency reporting evasion}: Cards are monetary instruments
				\item \textbf{Layering}: Purchased at multiple locations
				\item \textbf{Destination risk}: Flying to heightened scrutiny jurisdiction
			\end{itemize}
		\end{block}
	\end{frame}
	
	% ============================================================
	% SECTION 10: INSURANCE PRODUCT RISKS
	% ============================================================
	
	\begin{frame}{10.1: Life Insurance CDD Triggers}
		\begin{block}{When CDD Applies to Life Insurance}
			\begin{itemize}
				\item Annual premium exceeds USD/EUR 10,000
				\item Single premium exceeds USD/EUR 10,000
				\item Any suspicious transaction
			\end{itemize}
			\textbf{Who must be identified:}
			\begin{itemize}
				\item Policyholder
				\item Beneficiary (once designated)
				\item For legal entities: beneficial owners
			\end{itemize}
		\end{block}
		\textcolor{red}{Beneficiary CDD applies even if designated after policy issuance.}
	\end{frame}
	
	\begin{frame}{10.2: Life Insurance ML Exploitation Methods}
		\begin{block}{How Insurance Features Are Exploited}
			\textbf{Policy loan as layering:}
			\begin{itemize}
				\item Purchase policy with criminal proceeds
				\item Take policy loan for 90\% of cash value
				\item Loan proceeds are "legitimate" check
			\end{itemize}
			\textbf{Cash surrender extraction:}
			\begin{itemize}
				\item After brief holding period
				\item Surrender policy, receive check minus penalty
				\item Penalty is "laundering fee"
			\end{itemize}
		\end{block}
	\end{frame}
	
	\begin{frame}{10.3: Life Insurance AML Red Flags}
		\begin{block}{Life Insurance ML Indicators}
			\begin{itemize}
				\item Premium from offshore shell company
				\item Premium from secrecy haven jurisdiction
				\item Early surrender with proceeds to different account
				\item Policyholder is foreign PEP; beneficiary is shell company
				\item Large premium relative to legitimate income
				\item Premiums from multiple unrelated third parties
			\end{itemize}
		\end{block}
	\end{frame}
	
	\begin{frame}{10.4: Property/Casualty Insurance Risks}
		\begin{block}{Property/Casualty ML Risks}
			Fraudulent claims to convert illicit funds into legitimate claim payments.
		\end{block}
		\textbf{Key Vulnerabilities:}
		\begin{itemize}
			\item Insurance fraud generates proceeds that need laundering
			\item Over-valuation of claims
			\item Multiple claims for same loss
		\end{itemize}
	\end{frame}
	
	\begin{frame}{10.5: Reinsurance Risks}
		\begin{block}{Reinsurance ML Risks}
			Complex transactions that obscure the flow of funds across borders.
		\end{block}
		\textbf{Key Vulnerabilities:}
		\begin{itemize}
			\item Complex cross-border structures
			\item Multiple layers of reinsurance
			\item Difficulty tracing ultimate destination
		\end{itemize}
	\end{frame}
	
	\begin{frame}{10.6: Insurance CDD Requirements}
		\begin{block}{Insurance CDD Under FATF}
			Identification of the policyholder and beneficiary, with EDD for high-risk policies.
		\end{block}
		\textbf{Key Points:}
		\begin{itemize}
			\item Insurance companies are covered entities under FATF Rec. 10
			\item EDD for high-risk policies and PEPs
			\item Beneficiary identification is mandatory
		\end{itemize}
	\end{frame}
	
	% ============================================================
	% SECTION 11: VASP AND CRYPTOASSET RISKS
	% ============================================================
	
	\begin{frame}{11.1: VASP Definition}
		\begin{block}{Virtual Asset Service Provider (VASP)}
			A person or entity conducting virtual asset exchange, transfer, or custody services for third parties.
		\end{block}
		\textcolor{red}{FATF extends AML requirements to VASPs under \textbf{Recommendation 16}.}
	\end{frame}
	
	\begin{frame}{11.2: FATF Travel Rule for VASPs}
		\begin{block}{Virtual Asset Travel Rule (Recommendation 16)}
			\textbf{Required information for transfers above USD/EUR 1,000:}
			\begin{itemize}
				\item Originator: name, account/wallet address
				\item Originator: address, ID number, or DOB (at least ONE)
				\item Beneficiary: name and account/wallet address
			\end{itemize}
		\end{block}
		\textcolor{red}{\textbf{Receiving VASP obligations:}} Must obtain info; if not provided $\rightarrow$ follow up, delay, apply EDD.
	\end{frame}
	
	\begin{frame}{11.3: Cryptoasset ML Risks}
		\begin{block}{Cryptoasset Vulnerabilities}
			\begin{itemize}
				\item Pseudonymity and anonymity
				\item Rapid cross-border transfer
				\item Privacy coins and mixers obscure transactions
				\item DeFi protocols with no KYC
			\end{itemize}
		\end{block}
		\textcolor{red}{Privacy coins (Monero, Zcash) and mixers (Tornado Cash) are \textbf{challenges} for tracing.}
	\end{frame}
	
	\begin{frame}{11.4: Privacy Coins (Monero, Zcash)}
		\begin{block}{Anonymity-Enhanced Cryptocurrencies}
			\textbf{Monero (XMR):}
			\begin{itemize}
				\item Ring signatures - obscures which input is real
				\item Stealth addresses - one-time addresses per transaction
			\end{itemize}
			\textbf{Zcash (ZEC):}
			\begin{itemize}
				\item Zero-knowledge proofs (zk-SNARKs)
				\item Shielded transactions hide all details
			\end{itemize}
		\end{block}
	\end{frame}
	
	\begin{frame}{11.5: Crypto Mixers and Tumblers}
		\begin{block}{Mixing/Tumbling Process}
			\begin{enumerate}
				\item Customer sends crypto to mixing service
				\item Mixer pools funds from many users
				\item Mixer redistributes to new wallets controlled by customer
				\item Original source becomes obscured
			\end{enumerate}
		\end{block}
		\textbf{Red flags:} Use of mixer without legitimate explanation, refusal to explain source.
	\end{frame}
	
	\begin{frame}{11.6: NFT Wash Trading}
		\begin{block}{NFT Wash Trading Typology}
			\begin{enumerate}
				\item Self-controlled wallets buy/sell same NFTs among themselves
				\item Prices artificially escalate through circular transactions
				\item Creates fraudulent market record
				\item Sells NFT to unwitting third party at inflated price
				\item Converts to fiat as "capital gain"
			\end{enumerate}
		\end{block}
		\textbf{Exchange obligation:} File SAR based on wash trading pattern.
	\end{frame}
	
	\begin{frame}{11.7: DeFi and AML Responsibility}
		\begin{block}{FATF DeFi Guidance}
			If a person/entity maintains \textbf{control or sufficient influence} over a DeFi protocol:
			\begin{itemize}
				\item Through governance token majority (>50\%)
				\item Ability to upgrade contract code
				\item Ongoing management or control
			\end{itemize}
			\textcolor{red}{\textbf{They may be considered a VASP}} subject to AML obligations.
		\end{block}
	\end{frame}
	
	\begin{frame}{11.8: Cross-Chain Bridges and Layering}
		\begin{block}{Cross-Chain Bridge Exploitation}
			\begin{itemize}
				\item Criminal converts assets to token on privacy blockchain using bridge
				\item Privacy blockchain obscures transaction trail
				\item After mixing, funds bridged back as different asset
				\item Transactional link to original source is severed
			\end{itemize}
		\end{block}
		\textbf{Detection:} On-chain analysis can identify bridge usage patterns.
	\end{frame}
	
	\begin{frame}{11.9: Custodial vs. Non-Custodial Wallets}
		\begin{block}{AML Obligations by Wallet Type}
			\textbf{Custodial wallets:}
			\begin{itemize}
				\item VASP has full CDD obligations
				\item VASP controls funds and can monitor transactions
			\end{itemize}
			\textbf{Non-custodial wallets:}
			\begin{itemize}
				\item VASP's CDD ability is limited
				\item Cannot verify off-platform transactions
				\item Higher inherent AML risk
			\end{itemize}
		\end{block}
	\end{frame}
	
	\begin{frame}{11.10: Stablecoin Money Laundering Risks}
		\begin{block}{Stablecoin Vulnerabilities}
			\begin{itemize}
				\item Pegged to fiat currency - price stability facilitates value storage
				\item Rapid cross-border transfers without banking intermediaries
				\item Some issuers have weak or no CDD controls
				\item Can be redeemed for fiat at any time
			\end{itemize}
		\end{block}
		\textbf{Risk:} Cross-border value movement without traditional oversight.
	\end{frame}
	
	% ============================================================
	% SECTION 12: GAMING AND GAMBLING RISKS
	% ============================================================
	
	\begin{frame}{12.1: Casino AML Obligations (U.S. BSA)}
		\begin{block}{Casino Reporting Requirements}
			Casinos with revenue > \$1 million must:
			\begin{itemize}
				\item File CTRs for cash > \$10,000 in a gaming day
				\item File SARs for suspicious patterns > \$5,000
				\item Conduct CDD at \$3,000 chip purchase (FATF Rec. 12)
				\item Maintain AML compliance program
				\item Verify customer identity at threshold
			\end{itemize}
		\end{block}
	\end{frame}
	
	\begin{frame}{12.2: Chip Washing Scheme}
		\begin{block}{Chip Washing Pattern}
			\begin{itemize}
				\item Large cash chip purchase (e.g., \$300,000) - triggers CTR
				\item \textbf{Minimal gaming} (e.g., 45 minutes, \$1,500 wagered)
				\item Cash out remaining chips for casino check
				\item Accepts small "gaming loss" as laundering fee
			\end{itemize}
		\end{block}
		\textbf{Obligations:} File CTR AND file SAR.
	\end{frame}
	
	\begin{frame}{12.3: Online Gambling Risks}
		\begin{block}{Online Gambling ML Risks}
			Deposits and withdrawals via multiple payment methods, obscuring the source of funds.
		\end{block}
		\textbf{Key Vulnerabilities:}
		\begin{itemize}
			\item Cross-border payment complexity
			\item Multiple payment methods
			\item Small transactions below thresholds
		\end{itemize}
	\end{frame}
	
	\begin{frame}{12.4: Sports Betting Risks}
		\begin{block}{Sports Betting ML Risks}
			Match-fixing and insider betting generate proceeds that need laundering.
		\end{block}
		\textbf{Key Vulnerabilities:}
		\begin{itemize}
			\item Used to generate and launder funds
			\item Insider information abuse
			\item Multiple betting accounts
		\end{itemize}
	\end{frame}
	
	\begin{frame}{12.5: Gaming CDD Requirements}
		\begin{block}{Gaming CDD Requirements}
			Identification at set thresholds, transaction monitoring, and STR filing.
		\end{block}
		\textbf{Key Points:}
		\begin{itemize}
			\item Gaming establishments are covered entities under FATF
			\item CDD at chip purchase thresholds
			\item Transaction monitoring required
		\end{itemize}
	\end{frame}
	
	% ============================================================
	% SECTION 13: REAL ESTATE RISKS
	% ============================================================
	
	\begin{frame}{13.1: Real Estate ML Risks}
		\begin{block}{Real Estate Vulnerabilities}
			Purchase of property with illicit funds, over/under-valuation, and use of shell companies to hide ownership.
		\end{block}
		\textbf{Key Vulnerabilities:}
		\begin{itemize}
			\item Preferred integration vehicle
			\item All-cash purchases through anonymous LLCs
			\item Over/under-valuation techniques
		\end{itemize}
		\textcolor{red}{Real estate is a preferred \textbf{integration} vehicle.}
	\end{frame}
	
	\begin{frame}{13.2: Real Estate Shell Companies}
		\begin{block}{Shell Company Abuse}
			Shell companies conceal beneficial ownership and make tracing the source of funds difficult.
		\end{block}
		\textbf{Red Flags:}
		\begin{itemize}
			\item Recently formed LLC
			\item Registered agent is corporate service provider
			\item Refusal to disclose beneficial owner
			\item Multiple layers of ownership
		\end{itemize}
	\end{frame}
	
	\begin{frame}{13.3: FinCEN Geographic Targeting Orders (GTOs)}
		\begin{block}{GTO Requirements - Real Estate}
			\begin{itemize}
				\item Title companies must identify natural persons behind legal entities
				\item All-cash purchases above threshold (typically \$300,000)
				\item Targets specific metro areas (Miami, LA, NYC, etc.)
				\item \textbf{Temporary orders}, renewed periodically
			\end{itemize}
		\end{block}
		\textbf{Red flags:} Recently formed LLC, IOLTA wire, registered agent refuses UBO disclosure.
	\end{frame}
	
	\begin{frame}{13.4: Real Estate Over-Valuation}
		\begin{block}{Over-Valuation Technique}
			Property is purchased above market value, with the difference returned to the buyer as "clean" proceeds.
		\end{block}
		\textbf{Key Points:}
		\begin{itemize}
			\item Form of value transfer
			\item Appears as legitimate real estate transaction
			\item Difference is laundered proceeds
		\end{itemize}
	\end{frame}
	
	\begin{frame}{13.5: Real Estate Integration Example}
		\begin{block}{Integration Through Real Estate}
			\begin{enumerate}
				\item \textbf{Placement}: Structured cash deposits below \$10,000
				\item \textbf{Layering}: Funds into anonymous LLC
				\item \textbf{Layering/Integration}: LLC purchases property with cash
				\item \textbf{Integration}: Property generates rental income reported on taxes
				\item \textbf{Integration}: Property sold, proceeds deposited as legitimate
			\end{enumerate}
		\end{block}
	\end{frame}
	
	\begin{frame}{13.6: Real Estate EDD Requirements}
		\begin{block}{EDD for High-Risk Real Estate}
			Source of funds verification, beneficial ownership identification, and transaction monitoring.
		\end{block}
		\textbf{Key Points:}
		\begin{itemize}
			\item Real estate agents are covered under FATF Recommendation 22
			\item EDD for all-cash purchases
			\item Beneficial ownership identification required
		\end{itemize}
	\end{frame}
	
	% ============================================================
	% SECTION 14: GATEKEEPER RISKS
	% ============================================================
	
	\begin{frame}{14.1: Gatekeeper Definition}
		\begin{block}{What is a Gatekeeper?}
			Professionals who facilitate financial transactions or provide services that can be exploited for money laundering.
		\end{block}
		\textbf{Gatekeepers include:}
		\begin{itemize}
			\item Lawyers
			\item Accountants
			\item Trust and Company Service Providers (TCSPs)
			\item Notaries
		\end{itemize}
	\end{frame}
	
	\begin{frame}{14.2: Lawyer Gatekeeper Risks}
		\begin{block}{Lawyer Risks}
			Handling client funds in escrow accounts, creating structures to conceal ownership.
		\end{block}
		\textbf{Key Vulnerabilities:}
		\begin{itemize}
			\item IOLTA accounts are blind pass-through vehicles
			\item Professional status may bypass scrutiny
			\item Attorney-client privilege may be abused
		\end{itemize}
		\textcolor{red}{Lawyers are covered under \textbf{FATF Recommendation 22}.}
	\end{frame}
	
	\begin{frame}{14.3: IOLTA Account Vulnerabilities}
		\begin{block}{Why IOLTA Accounts Are Dangerous}
			\begin{itemize}
				\item Banks rely on attorney's professional status
				\item Do NOT look through pooled accounts
				\item Illicit funds enter credentialed by attorney status
				\item Multiple client funds commingled
				\item Rapid disbursement possible
			\end{itemize}
		\end{block}
	\end{frame}
	
	\begin{frame}{14.4: Notary Gatekeeper Risks}
		\begin{block}{Notary Risks}
			Authenticating documents that may be used to conceal or transfer ownership of illicit assets.
		\end{block}
		\textbf{Key Vulnerabilities:}
		\begin{itemize}
			\item Document authentication without scrutiny
			\item May facilitate shell company formation
			\item Covered under FATF Recommendation 22
		\end{itemize}
	\end{frame}
	
	\begin{frame}{14.5: Mitigating Gatekeeper Risks}
		\begin{block}{Risk Mitigation}
			Training, client acceptance committees, and standardized due diligence questionnaires.
		\end{block}
		\textbf{Key Controls:}
		\begin{itemize}
			\item AML training for professionals
			\item Client acceptance and review committees
			\item Standardized CDD questionnaires
			\item Ongoing monitoring
		\end{itemize}
	\end{frame}
	
	% ============================================================
	% SECTION 15: TRUST AND COMPANY SERVICE PROVIDER RISKS
	% ============================================================
	
	\begin{frame}{15.1: TCSP Definition}
		\begin{block}{Trust and Company Service Provider (TCSP)}
			An entity providing services such as company formation, registered office, nominee directors, and trust services.
		\end{block}
		\textcolor{red}{TCSPs are covered under \textbf{FATF Recommendation 23}.}
	\end{frame}
	
	\begin{frame}{15.2: TCSP ML Risks}
		\begin{block}{TCSP ML Risks}
			Formation of companies to conceal beneficial ownership and facilitate layering of illicit funds.
		\end{block}
		\textbf{Key Vulnerabilities:}
		\begin{itemize}
			\item Creates vehicles for money laundering
			\item Nominee directors and shareholders
			\item Registered office addresses
		\end{itemize}
	\end{frame}
	
	\begin{frame}{15.3: TCSP Red Flags}
		\begin{block}{Red Flags for TCSP Clients}
			\begin{itemize}
				\item Requesting confidentiality for shareholders
				\item High-risk jurisdictions
				\item Refusal to disclose beneficial ownership
				\item Complex structures without clear purpose
			\end{itemize}
		\end{block}
		\textcolor{red}{TCSPs should \textbf{decline} clients who refuse UBO disclosure.}
	\end{frame}
	
	\begin{frame}{15.4: TCSP CDD Requirements}
		\begin{block}{TCSP CDD Under FATF}
			Identify the client and beneficial owner, understand the purpose, and monitor the relationship.
		\end{block}
		\textbf{Key Points:}
		\begin{itemize}
			\item CDD is mandatory under FATF
			\item Identify beneficial owners
			\item Maintain records for 5+ years
			\item File SARs when suspicious
		\end{itemize}
	\end{frame}
	
	\begin{frame}{15.5: Trust Beneficial Ownership}
		\begin{block}{Trust CDD Requirements (FATF Rec. 25)}
			Must identify and verify:
			\begin{itemize}
				\item \textbf{Settlor}: Person who creates trust
				\item \textbf{Trustee}: Person/entity managing trust
				\item \textbf{Protector}: If has veto or replacement rights
				\item \textbf{Beneficiaries}: To extent determinable
			\end{itemize}
		\end{block}
		\textcolor{red}{Protector with veto qualifies as PEP if former head of state.}
	\end{frame}
	
	\begin{frame}{15.6: Bearer Shares - Highest Vulnerability}
		\begin{block}{Bearer Shares Risk}
			\begin{itemize}
				\item Transferred by physical delivery of certificate
				\item \textbf{NO public registry record} of ownership
				\item Ownership can be transferred anonymously
				\item FATF Rec. 24: prohibit or require immobilization
			\end{itemize}
		\end{block}
		\textbf{Bank obligations:}
		\begin{itemize}
			\item Verify custodian can identify holder
			\item Still need to identify UBO
			\item Decline account if custodian refuses disclosure
		\end{itemize}
	\end{frame}
	
	% ============================================================
	% SECTION 16: ACCOUNTANT RISKS
	% ============================================================
	
	\begin{frame}{16.1: Accountant Gatekeeper Risks}
		\begin{block}{Accountant ML Risks}
			Preparing accounts that conceal the source of funds and facilitating tax evasion.
		\end{block}
		\textbf{Key Vulnerabilities:}
		\begin{itemize}
			\item Can be used to conceal illicit activity
			\item Creating false records
			\item Overstating expenses
		\end{itemize}
	\end{frame}
	
	\begin{frame}{16.2: Accountant Red Flags}
		\begin{block}{Red Flags for Accountants}
			\begin{itemize}
				\item Clients who refuse to provide source of funds
				\item Complex structures without clear business purpose
				\item Refusal to provide documentation
				\item Unexplained wealth or income
			\end{itemize}
		\end{block}
		\textcolor{red}{Accountants must \textbf{question} unusual client requests.}
	\end{frame}
	
	\begin{frame}{16.3: Accountant ML Facilitation}
		\begin{block}{How Accountants Facilitate ML}
			\begin{itemize}
				\item Creating false records
				\item Overstating expenses
				\item Concealing beneficial ownership
				\item Preparing false tax returns
			\end{itemize}
		\end{block}
		\textbf{Key Points:}
		\begin{itemize}
			\item Accountants can be willing or unwilling facilitators
			\item Professional status may bypass scrutiny
			\item Covered under FATF Recommendation 22
		\end{itemize}
	\end{frame}
	
	\begin{frame}{16.4: Accountant Responsibilities}
		\begin{block}{Accountant AML Responsibilities}
			CDD, record-keeping, and STR filing for suspicious activity.
		\end{block}
		\textbf{Key Points:}
		\begin{itemize}
			\item Accountants are covered under FATF Recommendation 22
			\item Must conduct CDD on clients
			\item Must file SARs for suspicious activity
			\item Must maintain records
		\end{itemize}
	\end{frame}
	
	% ============================================================
	% SECTION 17: OTHER HIGH-RISK SECTORS AND STRUCTURES
	% ============================================================
	
	\begin{frame}{17.1: Free-Trade Zone Risks}
		\begin{block}{FTZ ML Risks}
			Lack of standardized inspection protocols for transshipped goods, enabling trade-based ML.
		\end{block}
		\textbf{Key Vulnerabilities:}
		\begin{itemize}
			\item Limited oversight
			\item Phantom shipping
			\item Cargo mislabeling
			\item Uniform invoice values
		\end{itemize}
		\textcolor{red}{FTZs are \textbf{high-risk} due to limited oversight.}
	\end{frame}
	
	\begin{frame}{17.2: Precious Metals and Stones Risks}
		\begin{block}{Precious Metals/Stones ML Risks}
			High-value, portable, and easily concealed assets that can be used for value transfer.
		\end{block}
		\textbf{Key Vulnerabilities:}
		\begin{itemize}
			\item High value to weight ratio
			\item Easily concealed
			\item Subjective valuation
		\end{itemize}
		\textcolor{red}{Precious metals are covered under \textbf{FATF Recommendation 23}.}
	\end{frame}
	
	\begin{frame}{17.3: Art and Antiquities Risks}
		\begin{block}{Art and Antiquities ML Risks}
			Subjective valuation makes over/under-invoicing easy, and sales can be anonymous.
		\end{block}
		\textbf{Key Vulnerabilities:}
		\begin{itemize}
			\item Value transfer and integration
			\item Subjective pricing
			\item Anonymous sales through auction houses
			\item Covered under FATF Recommendation 23
		\end{itemize}
	\end{frame}
	
	\begin{frame}{17.4: Shell Companies and Trusts}
		\begin{block}{Shell Company Risks}
			Conceal beneficial ownership and create opacity in financial transactions.
		\end{block}
		\textbf{Key Vulnerabilities:}
		\begin{itemize}
			\item Layering tool
			\item Bearer shares
			\item Nominee directors
			\item Multiple jurisdictions
		\end{itemize}
	\end{frame}
	
	\begin{frame}{17.5: Hawala and IVTS}
		\begin{block}{Hawala Characteristics}
			Informal value transfer system operating outside formal banking, based on trust and reciprocal obligations.
		\end{block}
		\textbf{Key Points:}
		\begin{itemize}
			\item No traceable banking record
			\item Operates outside formal financial system
			\item FATF Rec. 14 requires licensing/registration
		\end{itemize}
		\textcolor{red}{Hawala operators must register as MSBs - no cultural heritage exemption.}
	\end{frame}
	
	\begin{frame}{17.6: Money Mule Networks}
		\begin{block}{Money Mule Typology}
			A person who receives and transfers illicit funds through their account, often unknowingly recruited.
		\end{block}
		\textbf{Red Flags:}
		\begin{itemize}
			\item Multiple incoming P2P transfers from dozens of individuals
			\item Transfers typically under \$500 each
			\item Immediate aggregated wire transfers to offshore shell company
			\item Account holder with no verifiable income
		\end{itemize}
	\end{frame}
	
	% ============================================================
	% SUMMARY AND CONCLUSION SLIDES
	% ============================================================
	
	\begin{frame}{Domain A Summary: Top 30 Most Tested Concepts}
		\begin{multicols}{2}
			\begin{enumerate}
				\item ML vs. TF distinction
				\item Three stages: Placement, Layering, Integration
				\item TBML: Over/under-invoicing, phantom shipping
				\item BMPE uses drug dollars for legitimate invoices
				\item Nested correspondent accounts
				\item Shell banks prohibited
				\item PEPs: Family members included; SoW AND SoF required
				\item 6AMLD: 22 predicate offenses
				\item Travel Rule for VASPs
				\item DeFi: Control or influence = VASP obligations
				\item Privacy coins: Ring signatures, stealth addresses
				\item NFT wash trading
				\item Human trafficking: Multiple individuals to one controller
				\item Chip washing: Large buy-in, minimal gaming
				\item IOLTA vulnerabilities
				\item Bearer shares: Transfer by physical delivery
				\item SAR filing standard: Reasonable suspicion
				\item Tipping-off: Telling customer about SAR
				\item De-risking: Blanket termination violates RBA
				\item FATF effectiveness: IO.6, IO.7, IO.8
			\end{enumerate}
		\end{multicols}
	\end{frame}
	
	\begin{frame}{Critical Red Flags - Quick Reference}
		\begin{itemize}
			\item \textbf{Structuring:} \$9,000-\$9,800 repeated deposits
			\item \textbf{PEP:} Minister's daughter, "consulting fees," refuses SoW docs
			\item \textbf{TBML:} Price 400\% above/below market, shell company exporter
			\item \textbf{Human Trafficking:} Multiple individuals $\rightarrow$ 1 controller $\rightarrow$ multiple utility payments
			\item \textbf{NPO/TF:} Rapid cash withdrawal far from HQ
			\item \textbf{Correspondent:} Undisclosed nested transactions
			\item \textbf{DeFi:} 68\% governance token control, ability to upgrade code
			\item \textbf{Casino:} Large cash buy-in, minimal gaming, cash-out check
			\item \textbf{Real Estate:} Recently formed LLC, IOLTA wire, refuses UBO disclosure
			\item \textbf{Prepaid Cards:} Multiple cards loaded at \$9,500, border evasion
		\end{itemize}
	\end{frame}
	
	\begin{frame}{Exam Day Reminders for Domain A}
		\begin{itemize}
			\item \textbf{Domain A is 30\% of the exam} - largest single domain
			\item FATF Recommendations, PEPs, TBML, correspondent banking heavily tested
			\item \textbf{Select Two/Three/Four questions} = get ALL correct for credit
			\item Read scenarios carefully - facts matter
			\item Distinguish between ML and TF
			\item Know placement, layering, integration stages
			\item SARs filed on reasonable suspicion, not conclusive proof
			\item De-risking (blanket termination) is NOT preferred
			\item OSINT/adverse media does NOT require conviction
		\end{itemize}
		\vspace{0.5cm}
		\centering
		\textcolor{blue}{\Large Good luck on your CAMS exam!}
	\end{frame}
	
\end{document}